\documentclass[11pt,a4paper]{article}

% Fix fi/fl ligature extraction from the PDF (correct ToUnicode CMaps)
\input{glyphtounicode}
\pdfgentounicode=1
\usepackage{cmap}
\usepackage[utf8]{inputenc}
\usepackage[T1]{fontenc}
\usepackage{lmodern}
\usepackage{amsmath,amssymb,amsthm}
\usepackage{mathtools}
\usepackage{booktabs}
\usepackage{hyperref}
\usepackage{geometry}
\usepackage{enumitem}
\usepackage{listings}
\usepackage{longtable}
\usepackage{textcomp}
\usepackage{array}
\usepackage{xcolor}

\geometry{margin=1in}
\setcounter{tocdepth}{2}

\newtheorem{definition}{Definition}[section]
\newtheorem{principle}[definition]{Principle}
\newtheorem{theorem}[definition]{Theorem}
\newtheorem{proposition}[definition]{Proposition}
\newtheorem{invariant}[definition]{Invariant}
\newtheorem{remark}[definition]{Remark}
\newtheorem{condition}[definition]{Condition}

\lstset{
  basicstyle=\ttfamily\small,
  columns=fullflexible,
  keepspaces=true,
  xleftmargin=2em,
  aboveskip=0.5em,
  belowskip=0.5em,
  breaklines=true,
  showstringspaces=false,
  language=Haskell
}

\title{The Ledger: Addendum A1 to v10.3\\[6pt]
\large StatesHome --- Where Unit State Lives}
\author{R.\ Delloye}
\date{April 2026 --- Addendum A1 to v10.3}

\begin{document}
\maketitle

\begin{abstract}
Section~7 of the Ledger v10.3 specification leaves one question open: is the state of a
unit attached to the wallet or to the unit? Four instruments answer it by forcing distinct
keyings that no single attachment can carry --- a future whose state depends on past
trades, a managed account whose state is inherently per-client, a systematic strategy that
itself trades futures, and an instrument registered but never traded. Unit state therefore
lives in three maps, each with its own mutation discipline: an immutable, versioned
\emph{ProductTerms} keyed by unit; a mutable \emph{UnitStatus} keyed by unit and shared
across holders; and a \emph{PositionState} keyed by (wallet, unit). There is no
wallet-keyed state sector: every per-wallet economic fact is keyed by a mandate or strategy
unit and so lives in \emph{PositionState}. Twelve disciplines on these maps, collected in
\S\ref{sec:conditions}, render seven of the ten core invariants of v10.3 \S 11
unrepresentable rather than merely tested. This addendum supersedes the phrasing at v10.3
line 1034 and line 2287.
\end{abstract}

\tableofcontents
\clearpage

%==============================================================================
\section{The Question}
\label{sec:question}
%==============================================================================

\begin{quote}
Should the state of a unit be attached to the wallet or to the unit?
\end{quote}

\noindent Four instruments constrain the answer, and no two constrain it the same way:

\begin{enumerate}[leftmargin=*]
\item a future whose state depends on past transactions;
\item a managed account, whose state is inherently attached to a wallet;
\item a systematic (QIS) strategy that trades futures;
\item an instrument that exists in the universe but has not been traded yet.
\end{enumerate}

\noindent Each is worked in \S\ref{sec:testcases}. Each forces a part of the schema, and
the conditions that govern the schema are introduced where the instrument forces them, then
collected in \S\ref{sec:conditions}.

%==============================================================================
\section{Notation}
\label{sec:notation}
%==============================================================================

The framework primitives --- wallet, unit, holding, move, transaction, conservation ---
are those of v10.3. The symbols below are fixed for the remainder of this document.

\begin{longtable}{p{3.4cm}p{11.4cm}}
\toprule
\textbf{Symbol / term} & \textbf{Meaning} \\
\midrule
$\mathcal{W}$, $w$ & The set of wallets; a wallet. \\
$\mathcal{U}$, $u$ & The set of units; a unit. A \emph{unit} is any object that can be
held and conserved: an instrument, and also a mandate or a strategy contract
(\S\ref{sec:ma}, \S\ref{sec:qis}). \\
$h(w,u)$ & The signed quantity of unit $u$ held in wallet $w$, in exact minor units.
\emph{Conservation} is the framework invariant $\sum_{w\in\mathcal{W}} h(w,u)=0$ for every
$u$, with $h(w,u)=0$ for all but finitely many $w$. \\
\addlinespace
$\mathrm{Map}[K,V]$ & A finite partial function $K\rightharpoonup V$. Its accessor returns
$\mathrm{Option}[V]$. \\
$\mathrm{Option}[V]$ & $\mathrm{Some}\,v$ (a value $v$ is present) or $\mathrm{None}$ (no
value). \\
$\mathrm{NonEmptyList}[T]$ & A list of one or more elements of $T$. \\
total on $S$ & Defined for every key in $S$. \\
registration-total & Total on the set of \emph{registered} units, undefined elsewhere. \\
append-only & The only write is append; no key is overwritten in place and none is deleted. \\
monotone carrier & A map whose domain only grows: once a key is inserted it is never
removed. \\
\addlinespace
$\Delta f(w,u)$ & The change an event proposes to field $f$ of $\mathit{PositionState}[w,u]$;
$\Delta f(w,u)=0$ for every wallet the event does not touch. \\
conserved field & A \emph{PositionState} field that nets to zero across wallets:
\texttt{accumulated\_cost}, and \texttt{balance}. \texttt{balance} is a second conserved field,
moved by a transfer, carried only by the reference (\S\ref{sec:reference}) to exercise the C11
per-field-writer discipline with a canonical writer distinct from that of
\texttt{accumulated\_cost}; it is neither the framework holding $h(w,u)$ nor an economic datum
of the \S\ref{sec:answer} inventory. \texttt{hwm} and \texttt{entry\_nav} are non-conserved. \\
flat & A position whose conserved fields are zero. \\
$0_P$ & The \emph{PositionState} with every field zero (\texttt{zeroP},
\S\ref{sec:reference}). A first-touch row is $0_P$; a flat row need not be, since it may
retain a non-conserved field (a crystallised \texttt{hwm}, an \texttt{entry\_nav}). \\
$\mathit{StateDelta}$ & The change one event proposes, for a single unit, across the three
maps; conservation-checked and atomic (C2, C3). \\
\addlinespace
$u_{\mathrm{MA}}$ & A managed-account mandate, taken as a unit (\S\ref{sec:ma}). \\
$u_{\mathrm{QIS}}$ & A strategy contract, taken as a unit (\S\ref{sec:qis}). \\
$u_{\mathrm{ES}},u_{\mathrm{NQ}},u_{\mathrm{YM}}$ & Futures the strategy trades. \\
$u_{\mathrm{bench}}$ & A benchmark, taken as a unit. \\
$w_{\mathrm{QIS}}$, $w_C$ & The strategy's own wallet; a client wallet. \\
$u_{\mathrm{old}}\!\to\!u_{\mathrm{new}}$ & The pair an identity-breaking amendment
allocates (\S\ref{sec:untraded}). \\
\bottomrule
\end{longtable}

\noindent The holding symbol $h(w,u)$ is fixed here for this document; where v10.3 fixes a
symbol for quantity, that symbol governs.

%==============================================================================
\section{The Answer}
\label{sec:answer}
%==============================================================================

State attaches to three objects, each with its own map and its own mutation discipline.

\begin{lstlisting}[language={},basicstyle=\ttfamily\footnotesize]
ProductTerms  : Map[UnitId, NonEmptyList[TermsVersion]]   -- immutable, append-only, versioned; reg-total
UnitStatus    : Map[UnitId, UnitStatus]                   -- mutable, shared across holders; reg-total
PositionState : Map[(WalletId, UnitId), PositionState]    -- per (holder, unit); monotone, Option accessor
\end{lstlisting}

\noindent Each map's value type carries its sector's name by intent: a \emph{UnitStatus}
entry is one \texttt{UnitStatus} record, a \emph{PositionState} entry one
\texttt{PositionState} record. The reference names the maps themselves \texttt{ledgerUS} and
\texttt{ledgerPS} to keep the two apart (\S\ref{sec:reference}).

\noindent A fourth map carries identity, not state, and holds no financial value:

\begin{lstlisting}[language={}]
WalletRegistry : Map[WalletId, WalletMetadata]            KYC, permissions, audit cursor -- NOT state
\end{lstlisting}

\noindent The \emph{audit cursor} in \emph{WalletMetadata} is the position of the wallet's
reader in the event stream. \emph{WalletRegistry} carries no economic quantity and
participates in no conservation law.

There is no wallet-keyed state map. Call such a hypothetical map the \emph{$W$-sector}: a
$\mathrm{Map}[\mathrm{WalletId},\mathrm{WalletState}]$ holding per-wallet economic facts.
\S\ref{sec:ma} shows that every candidate datum for the $W$-sector is keyed by a mandate or
strategy unit, hence lives in \emph{PositionState}; the $W$-sector is empty of economic
state and is not built. The three maps are forced, and only three are forced, by
\S\ref{sec:why-three}.

\paragraph{Home of each datum.}
\begin{longtable}{p{3.4cm}p{2.6cm}p{8.4cm}}
\toprule
\textbf{Key} & \textbf{Home} & \textbf{Examples} \\
\midrule
$u$, immutable & \emph{ProductTerms} & multiplier, currency, expiry, clearinghouse, strike,
ISIN, fee schedule, mandate text, benchmark identity, index methodology \\
\addlinespace
$u$, mutable, shared & \emph{UnitStatus} & \texttt{lifecycle\_stage},
\texttt{last\_settlement\_price}, \texttt{last\_settlement\_date},
\texttt{current\_weights}, \texttt{nav\_index}, \texttt{triggered\_barrier},
\texttt{superseded\_by} \\
\addlinespace
$(w,u)$ & \emph{PositionState} & \texttt{accumulated\_cost}, \texttt{ccp\_binding},
per-position OTC lifecycle, \texttt{entry\_nav}, \texttt{hwm},
\texttt{accrued\_\{mgmt,perf\}\_fee}, \texttt{benchmark\_nav\_at\_inception},
\texttt{mandate\_breach\_flags} \\
\bottomrule
\end{longtable}

\noindent \texttt{balance} is absent from this inventory by intent: it is the reference's
demonstrative second conserved field (\S\ref{sec:notation}), exercising the per-field-writer
discipline, not an economic datum of the schema.

%==============================================================================
\section{The Four Instruments}
\label{sec:testcases}
%==============================================================================

Each condition is introduced at the instrument that forces it and stated once. The labels
C1--C12 are stable tags, not a sequence: they follow the \S\ref{sec:conditions} index rather
than order of first appearance, so the first met below is C2. They are load-bearing in
\S\ref{sec:unreachable} and \S\ref{sec:reference}. Parenthetical references to signals
S1--S4 point to the labelled expressibility notes in the reference (\S\ref{sec:reference}).

\subsection{A Future Whose State Depends on Past Transactions}
\label{sec:future}

Take the E-mini S\&P future. Its data divide by who may read them and how often they change.

\begin{longtable}{p{4.8cm}p{3.4cm}p{6.2cm}}
\toprule
\textbf{Field} & \textbf{Home} & \textbf{Reason} \\
\midrule
\texttt{multiplier}, \texttt{currency}, \texttt{expiry}, \texttt{clearinghouse},
\texttt{exchange}, \texttt{product\_id} & \emph{ProductTerms}$[u_{\mathrm{ES}}]$ &
Immutable. The CME contract and the ICE contract are \emph{distinct units}; the
clearinghouse is a term of the unit, not a per-wallet field (this replaces v10.3 line
1168). \\
\addlinespace
\texttt{lifecycle\_stage}, \texttt{last\_settlement\_price},
\texttt{last\_settlement\_date} & \emph{UnitStatus}$[u_{\mathrm{ES}}]$ & One settlement
price per contract, read identically by every holder. \\
\addlinespace
\texttt{accumulated\_cost} (\texttt{ac}) & \emph{PositionState}$[w,u_{\mathrm{ES}}]$ & Two
wallets holding the same contract carry different \texttt{ac}. \\
\bottomrule
\end{longtable}

\paragraph{Why per-(wallet, unit) is forced.}
Two wallets hold the same contract with different accumulated cost. Substituting one
wallet's position for the other's changes \texttt{ac}; any $u$-only keying identifies the
two and loses the difference. So \texttt{ac} is keyed by $(w,u)$, and \emph{PositionState}
exists.

\paragraph{Conservation of \texttt{ac}.}
\texttt{accumulated\_cost} is a conserved field (\S\ref{sec:notation}): conservation as
defined for $h$ extends to it, $\sum_{w\in\mathcal{W}}\texttt{ac}(w,u)=0$ for every $u$. No
map layout enforces this by types --- a refinement type on a sum of decimals is not free in
any production language. It is enforced one level up, at the event handler.

\begin{condition}[\textbf{C2}: handler-level conservation]
Every event handler --- \texttt{Trade}, \texttt{SettleVM}, \texttt{CorporateAction},
\texttt{QISRebalance}, \texttt{MandateAmend} --- emits a \texttt{StateDelta} whose every
conserved field satisfies
\[
\sum_{w\in\mathcal{W}} \Delta f(w,u) = 0
\]
structurally, by event class. The proof obligation per class covers the two-leg trade
($\Delta f = +q$ and $-q$), the $K$-leg trade ($K$ balanced legs), the variation-margin
fan-out (each wallet reset to target, cash legs offsetting), and the vacuous case of zero
touched wallets (the empty sum is $0$). Induction over the event stream lifts the per-event
identity to the global invariant.
\end{condition}

\noindent For the future: a \texttt{Trade} contributes buyer $\Delta\texttt{ac}=+q$ and
seller $\Delta\texttt{ac}=-q$, summing to $0$; a \texttt{SettleVM} resets each wallet to its
target with offsetting cash legs, summing to $0$. The worked legs are the proof.

\paragraph{What is not state.}
\texttt{first\_touch\_date} is derived from the event log on demand, not stored. Caching it
in \emph{PositionState} would make a replay disagree with itself under a back-dated
correction, because the cached value would reflect the pre-correction order while the fold
reflects the corrected one. The field is therefore absent from \emph{PositionState}.

\paragraph{Settled positions persist.}
A closed-out position keeps its row, flat, rather than being deleted. The row is read after
close-out for tax reporting, wash-sale lookback, and trade reconstruction. Two disciplines
on \emph{PositionState} make this precise; they are orthogonal --- one constrains the
accessor's type, the other constrains storage --- and both are required.

\begin{condition}[\textbf{C1}: Option accessor and monotone carrier]
\emph{(a) Option accessor.} \texttt{position(w,u)} returns
$\mathrm{Option}[\mathit{PositionState}]$. $\mathrm{None}$ means \emph{this wallet has never
held this unit}; $\mathrm{Some}(p)$ with $p$ flat means \emph{held once, now flat}. The two
readings are distinct and both load-bearing --- variation-margin settlement, wash-sale
lookback, and record-date entitlement each act on the difference --- so $\mathrm{None}$ and
$\mathrm{Some}(p)$ are never collapsed.\\
\emph{(b) Monotone carrier.} Once a row is created it is never removed; close-out leaves a
flat row. Replay is then a literal fold: \texttt{apply\_all(events)} is independent of where
checkpoints are cut, because the key set is stable across cuts, and conservation is a single
pass over that stable key set.
\end{condition}

\noindent The two halves are independent: (a) is a property of the return type, (b) a
property of the domain. Neither implies the other; the design carries both.

\paragraph{A canonical writer set per field.}
\texttt{ac} is written only by settle and trade handlers; a fee field is written only by the
fee handler; an entry-NAV field is written only at subscription. Concentrating each field's
writers removes cross-handler interference as a possibility rather than a thing to test.

\begin{condition}[\textbf{C11}: per-field canonical writer]
Each \emph{PositionState} field is tagged with its canonical writer set --- the closed set
of field-writers permitted to mutate it: \texttt{ac}$\to$settle/trade;
\texttt{balance}$\to$transfer; \texttt{hwm}$\to$fee\_crystallise;
\texttt{entry\_nav}$\to$subscribe. A write by any field-writer outside that set is a type
error at the writer's authorship site; the tag is erased once writes
share one delta row, so the guarantee binds at authorship, not at the stored row
(\S\ref{sec:reference}, signal S3). These field-writers (settle, trade, transfer,
fee\_crystallise, subscribe) are a different axis from the event classes of C2
(\texttt{Trade}, \texttt{SettleVM}, \texttt{CorporateAction}, \texttt{QISRebalance},
\texttt{MandateAmend}): one event class drives one or more field-writers, and the two
name-sets are not meant to coincide.
\end{condition}

\subsection{A Managed Account --- The ``Inherently Wallet-Attached'' Case}
\label{sec:ma}

The managed account is the case that appears to demand a $W$-sector: its state is the
client's, not the instrument's. It does not, because \emph{the mandate is a unit}. v10.3
\S 6 already places mandates in $\mathcal{U}$ (``the managed-account smart contract'', ``CSA
margin as a wallet-level smart contract''). The mandate $u_{\mathrm{MA}}$ is issued by the
manager and held by the client, and so conserves like any issued unit:
\[
h(w_{\mathrm{manager}},u_{\mathrm{MA}}) = -1,\qquad
h(w_{\mathrm{client}},u_{\mathrm{MA}}) = +1,\qquad
\sum_{w\in\mathcal{W}} h(w,u_{\mathrm{MA}}) = 0.
\]

\noindent With $u_{\mathrm{MA}}$ a unit, every field that looked per-wallet relocates to one
of the three maps:

\begin{longtable}{p{8.6cm}p{6.0cm}}
\toprule
\textbf{Field} & \textbf{Home} \\
\midrule
Mandate text, fee schedule, benchmark identity, max position limits &
\emph{ProductTerms}$[u_{\mathrm{MA}}]$ \\
HWM hurdle methodology, crystallisation frequency &
\emph{ProductTerms}$[u_{\mathrm{MA}}]$ \\
HWM \emph{value} (client-specific), \texttt{hwm\_date} &
\emph{PositionState}$[w_{\mathrm{client}},u_{\mathrm{MA}}]$ \\
Accrued mgmt / perf fee, mandate breach flags, subscription / redemption cursor &
\emph{PositionState}$[w_{\mathrm{client}},u_{\mathrm{MA}}]$ \\
Benchmark NAV at this wallet's inception &
\emph{PositionState}$[w_{\mathrm{client}},u_{\mathrm{MA}}]$ \\
Current benchmark level (shared, from index source) &
\emph{UnitStatus}$[u_{\mathrm{bench}}]$ \\
\bottomrule
\end{longtable}

\begin{condition}[\textbf{C12}: $W$-sector collapse]
All per-(wallet, mandate/strategy) economic state lives at
\emph{PositionState}$[w,u_{\mathrm{MA}}]$ or \emph{PositionState}$[w,u_{\mathrm{QIS}}]$.
No flat per-wallet scalar holds economic state; the empty $W$-sector is enforced by schema.
\end{condition}

\paragraph{The mandate-as-unit is not a sentinel.}
A rejected alternative (design C, \S\ref{sec:pareto}) inserts a sentinel unit with no
issuer, no holder, and no conservation partner, to carry wallet state. That sentinel breaks
$\sum_w h(\cdot)=0$ by fiat. $u_{\mathrm{MA}}$ is the opposite: a real contract with a real
issuer (the manager) and a real holder (the client), conserving by the standard issuance
law above. The collapse of the $W$-sector costs no conservation guarantee.

\paragraph{Two mandates, two rows.}
A client holding a base mandate and an overlay mandate on the same wallet carries two rows,
\emph{PositionState}$[w_{\mathrm{client}},u_{\mathrm{MA,base}}]$ and
\emph{PositionState}$[w_{\mathrm{client}},u_{\mathrm{MA,overlay}}]$, each with its own HWM,
fees, and breach flags. A flat per-wallet scalar would have identified them. This is the
direct reason C12 keys on the mandate unit and not on the wallet.

\subsection{A Strategy That Trades Futures}
\label{sec:qis}

A QIS strategy is itself a tradable unit, $u_{\mathrm{QIS}}$, and it also holds futures in
its own wallet. Five keyings coexist, and no two share a map:

\begin{lstlisting}[language={}]
u_QIS        the strategy itself (a tradable unit)
u_ES, u_NQ   the futures the strategy trades
u_MA,C       the mandate under which client C holds QIS (if any)
w_QIS        the strategy's own wallet (where its futures positions live)
w_C          a client wallet
\end{lstlisting}

\begin{longtable}{p{8.6cm}p{6.0cm}}
\toprule
\textbf{State} & \textbf{Home} \\
\midrule
Strategy terms (vol target, barrier, universe, share-class index start) &
\emph{ProductTerms}$[u_{\mathrm{QIS}}]$ \\
Strategy state (\texttt{last\_rebalance\_date}, \texttt{current\_weights},
\texttt{cumulative\_return}, \texttt{nav\_index}, \texttt{vol\_realised},
\texttt{triggered\_barrier}), shared across all holders &
\emph{UnitStatus}$[u_{\mathrm{QIS}}]$ \\
Futures the strategy holds (\texttt{accumulated\_cost}, ccp) &
\emph{PositionState}$[w_{\mathrm{QIS}},u_{\mathrm{ES}}]$, $[w_{\mathrm{QIS}},u_{\mathrm{NQ}}]$,
$[w_{\mathrm{QIS}},u_{\mathrm{YM}}]$ \\
Per-client subscription (\texttt{entry\_nav}, \texttt{hwm}) &
\emph{PositionState}$[w_C,u_{\mathrm{QIS}}]$ \\
Per-client mandate overlay (if any) &
\emph{PositionState}$[w_C,u_{\mathrm{MA,C}}]$ \\
\bottomrule
\end{longtable}

\paragraph{Where the HWM lives.}
A client holding two strategies carries two high-water marks. Any wallet-only keying
identifies them, exactly as in \S\ref{sec:ma}. The HWM therefore lives at
\emph{PositionState}$[w_C,u_{\mathrm{QIS}}]$, or
\emph{PositionState}$[w_C,u_{\mathrm{MA,C}}]$ when a mandate wrapper is in force.

\paragraph{The rebalance is one atomic event over several units.}
A barrier hit or monthly rebalance updates \emph{UnitStatus}$[u_{\mathrm{QIS}}]$ (new
weights, and a new lifecycle stage if the barrier breaks) \emph{and} the strategy's
\emph{PositionState}$[w_{\mathrm{QIS}},u_{\mathrm{ES}}]$,
$[w_{\mathrm{QIS}},u_{\mathrm{NQ}}]$, $[w_{\mathrm{QIS}},u_{\mathrm{YM}}]$ rows (the trades on
the underlying futures). These lie on different units; the event is therefore one
\texttt{StateDelta} per unit, applied together. A reader must never observe one applied and
another not.

\begin{condition}[\textbf{C3}: atomic StateDelta]
A \texttt{StateDelta} carries, for a single unit, the changes to \emph{ProductTerms},
\emph{UnitStatus}, and \emph{PositionState}. It applies in full or not at all; a partially
applied delta is not a representable state.
\end{condition}

\noindent \texttt{validate} discharges conservation for one unit across wallets
($\sum_{w}\Delta f(w,u)=0$). An event touching several units --- a rebalance trading
\texttt{ES}, \texttt{NQ}, and \texttt{YM}, or a future bought against cash --- is one
validated \texttt{StateDelta} per unit; per-unit conservation composes to event-level
conservation, and the group applies all-or-nothing as a single fold
(\S\ref{sec:reference}, signal S1).

\paragraph{Reads are scoped; exports flow through UnitStatus.}
A strategy publishes its state to clients through \emph{UnitStatus}$[u_{\mathrm{QIS}}]$, the
shared sector. A client's view of another client's overlay is not a read it is permitted to
make.

\begin{condition}[\textbf{C4}: capability-scoped reads]
Reads are capability-scoped. A cross-$(w,u_{\mathrm{MA}})$ overlay read is forbidden.
Strategy exports flow only through \emph{UnitStatus}.
\end{condition}

\paragraph{Wind-down.}
Setting \emph{UnitStatus}$[u_{\mathrm{QIS}}].\texttt{lifecycle\_stage}=\texttt{CLOSED}$
propagates to clients as a NAV-strike redemption. The
\emph{PositionState}$[w_C,u_{\mathrm{QIS}}]$ rows are retained, flat, by C1(b); they hold the
audit trail and the final HWM for tax reporting.

\subsection{An Instrument Registered but Not Yet Traded}
\label{sec:untraded}

An untraded option exercises the registration and lifecycle paths with no positions present.

\begin{itemize}[leftmargin=*]
\item \emph{ProductTerms}$[u]$ is present at registration.
\item \emph{UnitStatus}$[u]$ is present and fully initialised at registration with
product-declared defaults (e.g.\ \texttt{lifecycle\_stage = REGISTERED},
\texttt{last\_settlement\_price = None}); \texttt{unit\_status(u)} is total.
\texttt{REGISTERED} denotes a unit recorded in and known to the ledger before any
settlement --- an OTC option, a bespoke mandate, or an internal strategy reaches it
exactly as an exchange instrument does --- not admission to an exchange.
\item \emph{PositionState}$[w,u]$ has no rows; \texttt{position(w,u)} is $\mathrm{None}$ for
every $w$ until first touch (C1(a)).
\item \emph{WalletRegistry} is unaffected.
\end{itemize}

\begin{condition}[\textbf{C7}: ProductTerms registration-total]
\emph{ProductTerms} is registration-total: its first version is written at registration, so
\texttt{product\_terms(u)} is defined for every registered $u$.
\end{condition}

\begin{condition}[\textbf{C5}: UnitStatus registration-total]
\emph{UnitStatus} is registration-total: product-declared defaults are written at Unit Store
registration, so \texttt{unit\_status(u)} is defined for every registered $u$.
\end{condition}

\paragraph{Lifecycle with no holders.}
An untraded option moves $\texttt{REGISTERED}\to\texttt{ACTIVE}\to\texttt{EXPIRED}$ entirely
through \emph{UnitStatus}, creating no \emph{PositionState} row. The handlers fire on an
empty holder set and discharge conservation vacuously.

\begin{condition}[\textbf{C9}: vacuous conservation on zero-holder units]
A handler on a unit with no holders discharges $\sum_w\Delta f=0$ vacuously; the per-class
proof obligation of C2 includes the empty case. A handler that divides by holder count ---
the \texttt{dividend\_per\_share / len(holders)} bug class --- is excluded, because
conservation sums deltas and never divides by a count that may be zero.
\end{condition}

\paragraph{Registration is once only.}
\begin{condition}[\textbf{C10}: no re-registration]
Re-registering an existing unit id is a hard error, never a silent reset of its state.
\end{condition}

\paragraph{Amendment has two tracks.}
A bond amendment that preserves fungibility (coupon step-up, eligible-collateral change, a
fee tweak within a declared band) appends a new \texttt{TermsVersion} to
\emph{ProductTerms}$[u]$; existing \emph{PositionState} rows are untouched. An amendment that
breaks fungibility (new ISIN, benchmark swap, restructuring) allocates a fresh
$u_{\mathrm{new}}$ and stamps \texttt{superseded\_by}$(u_{\mathrm{old}}\to u_{\mathrm{new}})$
in \emph{UnitStatus}. Moving holders from $u_{\mathrm{old}}$ to $u_{\mathrm{new}}$ is a
\emph{separate} paired-issuance event, not part of the amendment: a burn on $u_{\mathrm{old}}$
and a mint on $u_{\mathrm{new}}$, each conserving on its own unit, applied together. A
single-unit \texttt{StateDelta} cannot span the two units (\S\ref{sec:reference}, signal S1).
The track is chosen by a product-declared total predicate,
\[
\texttt{is\_fungibility\_preserving} : \mathit{ProductTerms}\times\mathit{TermsAmendment}
\to \{\texttt{Preserving},\texttt{Breaking}\}.
\]

\begin{condition}[\textbf{C6}: ProductTerms append-only]
\emph{ProductTerms} is a $\mathrm{NonEmptyList}[\mathit{TermsVersion}]$, versioned and
append-only; no version is mutated in place.
\end{condition}

\begin{condition}[\textbf{C8}: two-track amendment]
The fungibility predicate decides the track: \texttt{Preserving} appends a
\texttt{TermsVersion} to the same unit; \texttt{Breaking} allocates a fresh $u$ and records
\texttt{superseded\_by}. The predicate is total and testable.
\end{condition}

%==============================================================================
\section{The Twelve Conditions}
\label{sec:conditions}
%==============================================================================

Each condition is defined once, at the instrument that forces it. This index gives the
label, its content in one line, and where it is established.

\begin{longtable}{p{1.2cm}p{10.6cm}p{1.8cm}}
\toprule
\textbf{} & \textbf{Content} & \textbf{Defined} \\
\midrule
C1 & Option accessor and monotone carrier, both. & \S\ref{sec:future} \\
C2 & Handler-level conservation, per event class, vacuous base case included. &
\S\ref{sec:future} \\
C3 & Atomic \texttt{StateDelta} across all three maps; partial application unrepresentable. &
\S\ref{sec:qis} \\
C4 & Capability-scoped reads; cross-overlay reads forbidden; exports through
\emph{UnitStatus}. & \S\ref{sec:qis} \\
C5 & \emph{UnitStatus} registration-total; product-declared defaults at registration. &
\S\ref{sec:untraded} \\
C6 & \emph{ProductTerms} versioned, append-only; no in-place mutation. &
\S\ref{sec:untraded} \\
C7 & \emph{ProductTerms} registration-total; first write at registration. &
\S\ref{sec:untraded} \\
C8 & Two-track amendment by a total fungibility predicate (Preserving / Breaking). &
\S\ref{sec:untraded} \\
C9 & Zero-holder handlers discharge $\sum=0$ vacuously; empty case in the proof obligation. &
\S\ref{sec:untraded} \\
C10 & Re-registration is a hard error, never a silent reset. & \S\ref{sec:untraded} \\
C11 & Each \emph{PositionState} field has a canonical writer set; a writer outside it is a
type error at authorship (erased at the row, S3). & \S\ref{sec:future} \\
C12 & All per-(wallet, mandate/strategy) economic state at
\emph{PositionState}$[w,u_{\mathrm{MA}}]$/$[w,u_{\mathrm{QIS}}]$; $W$-sector empty by schema. &
\S\ref{sec:ma} \\
\bottomrule
\end{longtable}

%==============================================================================
\section{Why Exactly Three Maps}
\label{sec:why-three}
%==============================================================================

Three independent constraints each break a distinct map; removing any one of the three
breaks the corresponding constraint.

\begin{enumerate}[leftmargin=*]
\item \textbf{Distinct per-holder state under the same contract.} Two wallets holding the
same contract carry distinct accumulated cost and lifecycle state (\S\ref{sec:future}).
$\Rightarrow$ a $(w,u)$-keyed map is required: \emph{PositionState}.
\item \textbf{Shared observables.} \texttt{last\_settlement\_price}, index values, and
strategy weights are one per contract and are read identically by every holder
(\S\ref{sec:future}, \S\ref{sec:qis}). $\Rightarrow$ a $u$-keyed map distinct from
per-holder state is required: \emph{UnitStatus}.
\item \textbf{Two mutation disciplines.} Terms are append-only and versioned for audit and
reconstruction; status is overwritten on every settle. Merging them puts two mutation
disciplines under one name (\S\ref{sec:untraded}). $\Rightarrow$ a third, append-only map is
required: \emph{ProductTerms}.
\end{enumerate}

\noindent A fourth map (an explicit $W$-sector) adds a sector whose economic content is
empty: every candidate datum is $(w,u_{\mathrm{MA}})$-keyed and collapses into
\emph{PositionState} (C12, \S\ref{sec:ma}). Three maps are the minimum basis of the problem,
not a compromise.

%==============================================================================
\section{Effect on v10.3, Lines 1034 and 2287}
\label{sec:line1034}
%==============================================================================

v10.3 line 1034 reads:

\begin{quote}
\textit{``Unit state is per-unit for most instrument types. For instruments with per-wallet
state (e.g., futures accumulated cost), the state dictionary is per (wallet, unit) pair.''}
\end{quote}

\noindent Replaced by:

\begin{quote}
Every unit $u\in\mathcal{U}$ carries immutable \emph{ProductTerms}$[u]$ (versioned,
append-only) and mutable \emph{UnitStatus}$[u]$ (shared across all holders). Every held
position $(w,u)$ carries \emph{PositionState}$[w,u]$, with $\mathrm{None}$ on the accessor
distinguishing ``never held'' from $\mathrm{Some}(0_P)$ ``held and flat''. State of a
holder's relationship to a strategy, mandate, or managed account lives at
\emph{PositionState}$[w,u_{\mathrm{MA}}]$, where $u_{\mathrm{MA}}$ is the mandate or strategy
unit. There is no separate wallet-keyed state sector.
\end{quote}

\noindent The migration is additive, then swap: \texttt{get\_unit\_state(u)} remains a
deprecated alias for the pair \texttt{(product\_terms(u), unit\_status(u))}; the
\texttt{get\_unit\_state(w,u)} of line 2287 maps to \texttt{position(w,u)}.

%==============================================================================
\section{Alternatives Considered and Rejected}
\label{sec:pareto}
%==============================================================================

The design is the fixed point of an adversarial multi-agent review (closing note). That
review ranked six designs on testability, correctness, and simplicity; the figures below are
its ordinal judgments on a 0--10 scale (higher better on all three) --- a relative ranking,
not measured quantities. Correctness is the gate axis. The per-design forcing reason that
follows the table, not the score, carries the argument.

\begin{longtable}{p{8.0cm}rrr}
\toprule
 & \textbf{Test.} & \textbf{Corr.} & \textbf{Simp.} \\
\midrule
A --- v10.3 current (per-unit, plus per-$(w,u)$ for futures only) & 4 & 5 & 4 \\
\textbf{B --- three maps + C1--C12 (adopted)} & \textbf{9} & \textbf{9} & \textbf{8} \\
C --- partial map with a sentinel unit & 7 & 3 & 7 \\
D --- four maps with an explicit $W$-sector & 7 & 7 & 5 \\
E --- state indexed only over the held set $\{(w,u): \texttt{position}(w,u)=\mathrm{Some}\}$ & 8 & 9 & 2 \\
F --- two maps with a universe wallet $w_\star$ & 5 & 6 & 6 \\
\bottomrule
\end{longtable}

\paragraph{The one forcing reason each rejected design loses.}
\begin{itemize}[leftmargin=*]
\item \textbf{A} has no sector for shared mutable observables distinct from immutable terms,
so settlement prices and immutable terms share one mutable cell --- the two mutation
disciplines of \S\ref{sec:why-three}(iii) collapse. Weaker than B on all three axes.
\item \textbf{C} (sentinel unit) fails day-one correctness: the sentinel has no issuer and
no conservation partner, so $\sum_w h(\cdot)=0$ fails by fiat (\S\ref{sec:ma}).
\item \textbf{D} (explicit $W$-sector) adds a sector whose economic content is empty: every
candidate datum is $(w,u_{\mathrm{MA}})$-keyed (\S\ref{sec:why-three}, C12).
\item \textbf{E} (state indexed only over the held set) equals B on correctness but its
construction has no implementation in available tooling, and is radically less simple.
Elegance does not decide; shippability does.
\item \textbf{F} (universe wallet $w_\star$) folds \emph{UnitStatus} into
\emph{PositionState}$[w_\star,u]$ and so re-introduces the per-unit / per-$(w,u)$ split as a
runtime predicate \texttt{is\_universe}$(w)$ on the wallet axis. Conservation
$\sum_w\Delta\texttt{ac}(w,u)=0$ must then exclude $w_\star$ explicitly --- a carve-out, not
a rule. \emph{UnitStatus} is the type-level expression of what F hides as a wallet-id
convention.
\end{itemize}

\noindent B strictly dominates A, C, D, and F on all three axes, and dominates E (equal
correctness, higher testability, far higher simplicity). B is the unique Pareto-optimum
under any correctness gate $\geq 7$, and is adopted.

%==============================================================================
\section{Invariants Made Unrepresentable}
\label{sec:unreachable}
%==============================================================================

Under design B with C1--C12, seven of the ten core invariants of v10.3 \S 11 are placed
beyond reach. The mechanism differs by invariant and is named in each gloss: an abstract type
whose only constructor is a checked smart constructor makes the illegal value unable to reach
the operation that would act on it (P1); a $\mathrm{NonEmptyList}$ makes the versionless state
untypable (P6); a per-field tag is a type error at authorship, then erased at the stored row
(P10). ``Unrepresentable'' is used in this precise sense, not as a claim that every guarantee
is a type-level fact: conservation, in particular, is a value-level check
(\S\ref{sec:reference}, signal S4). The invariant statements are those of v10.3 \S 11; the
gloss below names each, its mechanism, and the conditions that discharge it.

\begin{description}[leftmargin=3em,labelwidth=2em,labelsep=0.5em,style=nextline]
\item[\textbf{P1}] (conservation of quantity) --- handler-level per-class zero-sum (C2) and
atomic delta (C3). \texttt{ValidDelta} is abstract; its only constructor, \texttt{validate},
discharges the zero-sum, so an unconserved delta cannot reach \texttt{applyDelta}. The check
is value-level (signal S4), not a type fact.
\item[\textbf{P3}] (determinism of replay) --- replay is the \texttt{foldM} of the event
stream, hence a fold homomorphism: \texttt{replay (xs <> ys) = replay xs >=> replay ys},
where \texttt{f >=> g} is the composition that runs the error-returning step \texttt{f},
feeds its result to \texttt{g}, and stops at the first error (composition of
\texttt{Either LedgerError}-returning steps). Replaying a concatenated log equals replaying
each part and composing the two. The monotone carrier (C1(b)) keeps the key set stable
across any checkpoint cut, so checkpoint-independence is a consequence of this law, not a
test.
\item[\textbf{P5}] (idempotency of lifecycle events) --- the lifecycle stage lives in
\emph{UnitStatus}$[u]$, $u$-keyed and written by replacement, so re-applying a stage write is
the identity ($\texttt{EXPIRED}$ over $\texttt{EXPIRED}$ is $\texttt{EXPIRED}$): idempotency is
structural at a single key, not cross-map coordination, and an untraded unit's lifecycle
(\S\ref{sec:untraded}) discharges it touching no \emph{PositionState} row. The non-conserved
\emph{PositionState} fields carry the same algebra --- \texttt{hwm} by $\max$
($\max(x,x)=x$), \texttt{entry\_nav} write-once --- while the additive conserved fields
\texttt{accumulated\_cost} and \texttt{balance} draw their replay-safety from P1, not from
replacement.
\item[\textbf{P6}] (immutability of terms) --- append-only
$\mathrm{NonEmptyList}[\mathit{TermsVersion}]$ (C6) and registration totality (C7).
\item[\textbf{P7}] (no in-place mutation of identity) --- re-registration rejected (C10);
breaking amendments allocate $u_{\mathrm{new}}$ and never rewrite $u_{\mathrm{old}}$ (C8).
\item[\textbf{P9}] (capability scoping) --- cross-$(w,u_{\mathrm{MA}})$ overlay reads
forbidden (C4); enforced at the capability layer wrapping the accessors, not in the data
reference (signal S2).
\item[\textbf{P10}] (handler-field canon) --- the canonical writer set per field (C11); a
write by a field-writer outside it is a type error at authorship, erased once writes share a
delta row (signal S3).
\end{description}

\noindent No other candidate design makes more than three of the ten unrepresentable.

%==============================================================================
\section{Testing Commitments}
\label{sec:mutationscore}
%==============================================================================

The specification commits to these figures.

\begin{itemize}[leftmargin=*]
\item Event handlers (arithmetic on exact decimals): mutation score \textbf{85--90\%} ---
conservation kills sign and coefficient mutants.
\item Lifecycle guards ($\texttt{REGISTERED}\to\texttt{ACTIVE}\to\texttt{EXPIRED}$):
\textbf{70--80\%} --- boundary-comparison mutants (\texttt{>} vs \texttt{>=}) require
targeted tests. The lifecycle is a flat enumeration, so transition ordering (no
$\texttt{EXPIRED}\to\texttt{REGISTERED}$) is enforced by these tests, not by types --- distinct
from P5, whose idempotency is structural.
\item Core overall: $\geq$\textbf{80\%} (Feathers threshold).
\item State-machine model at $|W|=3,\ |U|=2,\ \text{depth}\leq 6$: $10^5$--$10^6$ reachable
states (TLC-tractable); conservation-violating handlers caught at depth 1.
\end{itemize}

\noindent The generator universe (CDM enum $\times$ product-type) is finite and enumerable,
so coverage is structurally bounded.

%==============================================================================
\section{Risk Register}
\label{sec:risk}
%==============================================================================

Eight migration, governance, and operational risks. None is a correctness blocker; each is
budgeted.

\begin{longtable}{p{0.7cm}p{6.3cm}p{7.5cm}}
\toprule
\textbf{\#} & \textbf{Risk} & \textbf{Mitigation} \\
\midrule
F1 & Migration scope: every downstream reader of \texttt{unit\_state}. & Ship
\texttt{get\_unit\_state} as a deprecated alias; split additively before cutover;
parallel-run at least one quarter with reconciliation. \\
\addlinespace
F2 & Ownership of the C8 fungibility predicate is ungoverned. & Assign per-product RACI in
the Unit Store registry; require Legal/Product/Risk sign-off for predicate changes; version
the predicate. \\
\addlinespace
F3 & Monotone-carrier key-space growth at $10^7$ wallets $\times\,10^6$ units is
unbenchmarked. & Benchmark Option-accessor cost and cache behaviour before merge; add a
tombstone-compaction policy for rows held at $\mathrm{Some}(0_P)$ beyond the retention
horizon. \\
\addlinespace
F4 & C1--C12 is a multi-sprint programme, not one release. & Land the schema with
C1/C5/C6/C7/C9/C10 (type-enforced); stage C2/C3/C8/C11 over later releases; never merge
without at least C1, C2, C3. \\
\addlinespace
F5 & Mandate-as-unit creates an SFTR / EMIR reporting surface for mandate issuance. &
Pre-flight with the Regulatory team; determine whether $u_{\mathrm{MA}}$ issuance requires a
UTI / LEI pair; consider a \texttt{reportable} flag on
\emph{ProductTerms}$[u_{\mathrm{MA}}]$. \\
\addlinespace
F6 & CDM alignment (per-trade \texttt{TradeState} vs \emph{PositionState}$[w,u]$) is
asserted, not verified. & Re-run the CDM mapping (Rosetta NS1--7) against the three-map schema
before any CDM adapter work; publish a delta document. \\
\addlinespace
F7 & Onboarding cost of three maps, Option/monotone, and twelve conditions is real. & Ship a
decision tree and runbook; require training before handler work; provide a one-page cheat
sheet. \\
\addlinespace
F8 & Forward migration is irreversible once the $W$-sector is collapsed. & Keep a read-only
mirror of the prior wallet-state for at least four quarters post-cutover; document the
inverse mapping (overlay key $u_{\mathrm{MA}}\to$ former field) explicitly. \\
\bottomrule
\end{longtable}

\noindent F2, F5, and F6 require external stakeholders before any code ships. F1, F3, F4,
F7, and F8 are internal engineering risks, manageable within the delivery programme.

%==============================================================================
\section{Reference Implementation}
\label{sec:reference}
%==============================================================================

The reference is \texttt{reference/StatesHome.hs}, included below. It is total and
deterministic, and makes the illegal states unrepresentable rather than merely checked.
Quantities are exact integer minor units, not floating point, so determinism and
conservation are arithmetic facts.

The encoding carries the conditions structurally:

\begin{itemize}[leftmargin=*]
\item \texttt{Ledger} is abstract with no row deleter, so no caller can remove a
\emph{PositionState} row --- the monotone half of C1(b). \texttt{position} returns
\texttt{Maybe} --- the Option half of C1(a).
\item \texttt{ValidDelta} is abstract; its only constructor is \texttt{validate}, which
discharges conservation as a monoid identity (C2). An unconserved delta cannot reach
\texttt{applyDelta}. The empty fold is \texttt{mempty}, so zero-holder events validate with
no special case (C9).
\item \texttt{ProductTerms} is abstract over a \texttt{NonEmpty}; the only growth operations
are \texttt{register} (singleton, C7) and \texttt{appendVersion} (append, C6). No setter
exists.
\item The \texttt{FieldWrite h} GADT constructors \emph{are} the C11 field-to-writer table: a
write by the wrong handler fails to typecheck at the handler's authorship site.
\texttt{settleHandler} has output type \texttt{Map WalletId [FieldWrite 'Settle]} and is the live
call site that exercises this: \texttt{main} builds \texttt{tradeSD} and \texttt{closeSD} as
\texttt{erase\,.\,settleHandler}, erasing the index through
\texttt{erase\,=\,fmap\,(map\,SomeWrite)} as the writes enter \texttt{sdRows}, so the guarantee
binds at authorship, not at the delta row (signal S3). The commented \texttt{\_c11\_bad} witnesses
the rejected cross-handler write.
\item \texttt{applyDelta} returns the whole new ledger or a \texttt{Left}; there is no
partial state (C3). \texttt{register} rejects re-registration with \texttt{ReRegistration}
(C10) and writes \emph{ProductTerms} and \emph{UnitStatus} together, so ``registered in PT
$\iff$ registered in US'' holds by construction.
\item \texttt{amend} is the two-track C8: \texttt{Preserving} appends; \texttt{Breaking}
allocates a fresh unit and stamps \texttt{usSupersededBy}, never rewriting the old terms.
\end{itemize}

\lstinputlisting[language=Haskell]{reference/StatesHome.hs}

%==============================================================================
\section{The One-Sentence Answer}
\label{sec:onesentence}
%==============================================================================

\begin{quote}
\textbf{Unit state lives in three places --- immutable \emph{ProductTerms}$[u]$, shared
mutable \emph{UnitStatus}$[u]$, and per-position \emph{PositionState}$[w,u]$ --- with no
separate wallet-keyed state sector, because every per-wallet economic fact is keyed by a
mandate or strategy unit and so lives in \emph{PositionState}$[w,u_{\mathrm{MA}}]$.}
\end{quote}

\noindent This is the Pareto-optimal schema on testability, correctness, and simplicity, and
the minimum basis of the problem, not a compromise.

\vfill
\begin{center}
\rule{0.5\textwidth}{0.4pt}\\[4pt]
\textit{The design is the fixed point of adversarial multi-agent review; the supporting
analyses are archived under \texttt{StatesHome\_work/}.}
\end{center}

\end{document}
